%------------------------------------------------------------------
% Asymptotic behavior of the Lerch transcendent function,
% by
% L. M. Navas, F. J. Ruiz and J. L. Varona
%
% Date of first version: January 11, 2011
% Date of last modification: February 10, 2012
%
%------------------------------------------------------------------


\documentclass[11pt,a4paper,reqno]{amsart}

\usepackage{url}
%\usepackage{hyperref}
\usepackage[colorlinks,linkcolor=blue,citecolor=blue]{hyperref}


%\date{February 10, 2012}
\renewcommand{\datename}{\textbf{THIS PAPER HAS BEEN PUBLISHED IN:}}
\date{\textit{J.\@ Approx.\@ Theory} \textbf{170} (2013), 21--31.}

%--------------------------------------
%% ONLY FOR DRAFTS:
%\usepackage[notcite,notref]{showkeys}
%--------------------------------------

%--------------------------------------
%% ONLY FOR THE SPANISH DRAFTS:
%\usepackage[latin1]{inputenc}
%\usepackage[spanish]{babel}
%--------------------------------------


%--------------------------------------
%% ONLY FOR DRAFTS:
%\usepackage[svgnames]{xcolor}
%\newcommand{\repasar}[1]{\textcolor{red}{#1}}
%\newcommand{\comentario}[1]{\textcolor{Green}{#1}}
%\newcommand{\quitar}[1]{\textcolor{blue}{#1}}
%\newcommand{\nuevo}[1]{\color{DarkViolet}{#1}\color{black}}
%--------------------------------------


\newcommand{\N}{\mathbb{N}}
\newcommand{\Z}{\mathbb{Z}}
%\newcommand{\Q}{\mathbb{Q}}
\newcommand{\R}{\mathbb{R}}
\newcommand{\C}{\mathbb{C}}
\newcommand{\dom}{\mathbb{C}} % Domain (to change it if necessary)

\newcommand{\B}{\mathcal{B}}
\newcommand{\E}{\mathcal{E}}

%\renewcommand{\epsilon}{\varepsilon}

\renewcommand{\Re}{\operatorname{Re}}
\renewcommand{\Im}{\operatorname{Im}}
%\newcommand{\re}{\operatorname{Re}}
%\newcommand{\im}{\operatorname{Im}}
\newcommand{\Li}{\operatorname{Li}}

\theoremstyle{plain}
\newtheorem{theorem}{Theorem}%[section]
%\newtheorem{prop}[theorem]{Proposition}
\newtheorem{corollary}[theorem]{Corollary}
\newtheorem{lemma}{Lemma}

\theoremstyle{remark}
\newtheorem{remark}{Remark}


%--------------------------------------
\title[Asymptotic behavior of the Lerch transcendent function]%
{Asymptotic behavior of the Lerch transcendent function}

%\author{L.~M.~Navas, F.~J.~Ruiz and J.~L.~Varona}

\author[L.~M.~Navas]{Luis M. Navas}
\address{%
        Departamento de Matem\'aticas,
        Universidad de Salamanca,
        Plaza de la Merced 1-4,
        37008 Salamanca, Spain}
\email{navas@usal.es}

\author[F.~J.~Ruiz]{Francisco J. Ruiz}
\address{%
        Departamento de Matem\'aticas,
        Universidad de Zaragoza,
        Campus de la Plaza de San Francisco,
        50009 Zaragoza, Spain}
\email{fjruiz@unizar.es}

\author[J.~L.~Varona]{Juan L. Varona}
\address{%
        Departamento de Matem\'aticas y Computaci\'on,
        Universidad de La Rioja,
        Calle Luis de Ulloa s/n,
        26004 Logro\~no, Spain}
\email{jvarona@unirioja.es}
\urladdr{http://www.unirioja.es/cu/jvarona/welcome.html}

\thanks{The research of the second and third authors was supported by grant MTM2009-12740-C03-03 of the~DGI}

\keywords{Lerch transcendent function, asymptotic estimates,
Fourier series}

\subjclass[2010]{Primary 41A60; Secondary 11M35, 42A10}
%--------------------------------------

%--------------------------------------
\begin{document}
%--------------------------------------

%--------------------------------------
\begin{abstract}
    For complex parameters $\lambda$ and $s$, consider the Lerch transcendent
  $\Phi(\lambda,s,z) = \sum_{k=0}^\infty \lambda^k (k+z)^{-s}$
  as a function of the complex variable~$z$.
  We analyze the asymptotic behavior of this function as
  $\operatorname{Re} s \to -\infty$.
\end{abstract}
%--------------------------------------

%--------------------------------------
\maketitle
%--------------------------------------

\section{Introduction}

The Lerch transcendent is given by the series
\begin{equation}
\label{eq:Lerch}
  \Phi(\lambda,s,z) = \sum_{k=0}^\infty\frac{\lambda^k}{(k+z)^s};
\end{equation}
see \cite[\S1.11, p.~27]{Bate} or \cite[\S25.14]{Ap-NIST}, for example.
This function, defined by Mathias Lerch in 1887 in his paper \cite{Le},
includes as special cases of the parameters the Hurwitz and Riemann zeta functions
and the polylogarithms, among others, and therefore
has applications ranging from number theory to physics.
It is often used to obtain new identities; see, for instance, \cite{BaSa, Boy, GuiS, SC-book}.
Several kinds of asymptotic estimates for $\Phi$ have been studied; see~\cite{FL} and the references therein.

Here we shall consider $\Phi$ as a function of the complex variable $z$, with
$\lambda$ and $s$ as parameters. We shall assume that the summands
\[
  \frac{1}{(k+z)^s} = e^{-s\log(k+z)},
  \quad k\in\N\cup\{0\},
\]
are evaluated using the principal branch of the logarithm, and thus are holomorphic
on $\C\setminus({-\infty}, 0]$.

We are interested in the behavior as $|s| \to \infty$ in a region with bounded imaginary part and with real part $\Re s \to -\infty$.
Since
\begin{equation}
\label{eq:sreal}
  \left|(k+z)^{-s}\right| = |k+z|^{-\Re s}e^{\Im s \arg(z+k)},
\end{equation}
the fact that $\Im s$ remains bounded means that this modulus is controlled
by $\Re s$. This restriction will often allow us to reduce to the case of negative real~$s$ when obtaining estimates.

As far as the parameter $\lambda$ is concerned, since $\Phi(0,s,z) = z^{-s}$, we may
exclude $\lambda = 0$ as a trivial case. If $|\lambda| < 1$ then, for any $s \in \C$, the series~\eqref{eq:Lerch} converges uniformly in $z$ on compact subsets of $\C\setminus ({-\infty},0]$, and hence defines a holomorphic function in this region.
We may enlarge the domain to allow $z \in \C\setminus \{0,-1,-2,\dots\}$ by including the branch discontinuity of the principal argument, at the cost of losing holomorphy.
This is standard in the mathematical literature, see for example the above mentioned~\cite{Ap-NIST,Bate}.

In addition, for $\Re s < 0$, the summand $(k+z)^{-s}$ in~\eqref{eq:Lerch} can be continuously extended to $z = -k$ by defining $0^{-s} = 0$.
In what follows, when $\Re s < 0$, we shall assume that $\Phi(\lambda,s,z)$ has been defined for all $z \in \C$ in this way.

Given the previous remarks, one sees that, for $|\lambda| < 1$ and $\Re s < 0$,
$\Phi(\lambda,s,x)$ extends to a continuous function of $x \in [0,1]$ and, in particular, it is absolutely integrable. In fact, absolute integrability holds for
$\Re s < 1$, since, although we lose continuity at $x = 0$, we still have
\[
  \int_0^1 |\Phi(\lambda,s,x)|\,dx
  \leq \sum_{k=0}^\infty |\lambda|^k \int_0^1 \frac{dx}{(k+x)^{\Re s}} < +\infty
\]
and for fixed $a > 0$, the series $\sum_k\lambda^k\,k^a$, considered as a power series in $\lambda$, has radius of convergence equal to~$1$.

The paper is structured as follows. In Section~\ref{sec:SF}, we find the Fourier series of $\lambda^x \Phi(\lambda,s,x)$ for $x\in[0,1]$, which is essentially the Lindel\"of-Wirtinger expansion of $\Phi$ (Theorem~\ref{th:serieFourier}); this approach follows~\cite{NRV-JAT, NRV-MC} and gives an alternative to the classical complex analytic methods as found for instance in~\cite[formula 1.11\,(6), p.~28]{Bate}. It is then shown that the Fourier series gives an asymptotic expansion for
$\Phi(\lambda,s,z)$ as $s \to -\infty$, valid for $z = x\in[0,1]$ (see Corollary~\ref{cor:FS-admissible}). In Section~\ref{sec:extension}, we prove the validity of this asymptotic expansion for arbitrary complex-valued $z$, and we show that it is uniform on compact subsets of~$\C$; this is expressed in Theorem~\ref{th:FS-admissible}, which is the main result of the paper. Finally, in Sections~\ref{sec:asymptotic} and~\ref{sec:polylog}, we consider some consequences of the main theorem and its application to special cases of the Lerch transcendent, such as the polylogarithm functions.


\section{The Fourier series and asymptotics of $\Phi$ on $[0,1]$}
\label{sec:SF}


Let us begin by noting that Fourier analysis provides an easy and direct way of obtaining the Lindel\"of-Wirtinger expansion (Theorem \ref{th:serieFourier}) for the Lerch transcendent. This important result is the starting point of our study. It can be found for instance in \cite[p.~28]{Bate} or \cite[p.~34]{MOS}, and a special case already appears in ~\cite{Wir}. For completeness and in order to focus attention on the use of Fourier analytic methods, we include a brief overview of the main results, detailing precisely the domains of validity.


We only need a few simple facts, the first of which is the well-known formula
%
%\begin{lemma}
%\label{lem:integral}
% If % <-- al quitar el "Lema" lo cambio por "for"
\begin{equation}
\label{eq:integral}
  \int_0^\infty e^{-ta}\,t^{-s}\,dt = a^{s-1}\Gamma(1-s),
  \qquad
  \Re s < 0, \ \Re a > 0
\end{equation}
%\end{lemma}
(this is \cite[formula 1.5.1\,(33), p.~12]{Bate}); for $a > 0$ it is immediate by changing variables in the definition of $\Gamma(1-s)$, and follows for $\Re a > 0$ in general by analytic continuation.
% \begin{proof}
% If $a \in \R$, $a>0$, simply change variables in
% $\Gamma(1-s) = \int_{0}^{\infty} t^{-s} \, e^{-t} \, dt$.
% The general case of $\Re a>0$ follows by analytic continuation.
% \end{proof}
With this, obtaining the Fourier coefficients of $\lambda^{x} \Phi(\lambda,s,x)$ is an immediate computation.

\begin{lemma}
Let $\lambda$ and $s$ be complex parameters with $0<|\lambda|<1$ and $\Re s<0$.
For every $n\in\Z$, we have
\[
  \int_0^1 \lambda^x\,\Phi(\lambda,s,x)\,e^{-2\pi inx}\,dx
  = (2n\pi i-\log\lambda)^{s-1} \Gamma(1-s).
\]
\end{lemma}

\begin{proof}
Using~\eqref{eq:Lerch} and~\eqref{eq:integral}, we have
\begin{multline*}
  \int_0^1 \lambda^x\,\Phi(\lambda,s,x)\,e^{-2\pi inx}\,dx
  = \sum_{k=0}^\infty \int_0^1 \frac{\lambda^{k+x}}{(k+x)^s}e^{-2\pi in(k+x)}\,dx
  \\
  =\int_0^\infty \lambda^t\,e^{-2\pi int}t^{-s}\,dt
  = \int_0^\infty e^{-t(2\pi in-\log\lambda)}t^{-s}\,dt
\end{multline*}
and applying \eqref{eq:integral} finishes the proof.
One easily verifies that the exchange of sum and integral is justified.
Note that we need $\log|\lambda|<0$, which is indeed the case.
\end{proof}

Finally, our hypotheses on the parameters imply that
\[
  \Phi(\lambda,s,0) = \sum_{k=1}^\infty \frac{\lambda^k}{k^s} = \lambda \Phi(\lambda,s,1),
\]
hence $\lambda^x \Phi(\lambda,s,x)$ extends to a continuous function on $[0,1]$, and since it is of class $\mathcal{C}^{(1)}$ on $(0,1)$, by Dirichlet's basic theorem on the convergence of Fourier series, we obtain:

\begin{theorem}[Lindel\"of-Wirtinger expansion]
\label{th:serieFourier}
Let $\lambda$ and $s$ be complex parameters with $0<|\lambda|<1$ and $\Re s<0$. Then
\[
  \Phi(\lambda,s,x)
  = \lambda^{-x}\,\Gamma(1-s) \sum_{n\in\Z} (2n\pi i-\log\lambda)^{s-1}e^{2n\pi i x},
  \quad x\in[0,1],
\]
uniformly on $[0,1]$.
\end{theorem}

\begin{remark}
When $s=1-k$ with $k\in\N$, the series defining $\Phi(\lambda,s,z)$
actually sums to a polynomial in $z$ of degree $k-1$, with coefficients that are rational functions in $\lambda$ with a unique pole at $\lambda=1$. In this case,
$\Phi(\lambda,s,z)$ extends to all values of $\lambda\neq 0,1$. The relation
\[
  \Phi(\lambda,1-k,z) = -\frac{\B_k(z;\lambda)}{k}
\]
defines the Apostol-Bernoulli polynomials $\B_{k}(z;\lambda)$
(see~\cite{Ap}). Thus we obtain, in particular, for $k \geq 2$,
\[
  \B_k(x;\lambda) = -\lambda^{-x}\,k! \sum_{n\in\Z} \frac{e^{2n\pi i x}}{(2n\pi i-\log\lambda)^k},
  \quad x\in[0,1],
\]
under the restriction $|\lambda|<1$. In fact the formula is valid for all
$\lambda\neq 0,1$, as was shown in~\cite{NRV-MC} using certain algebraic properties of this polynomial family.
\end{remark}

\begin{remark} The Fourier series in Theorem~\ref{th:serieFourier} has two alternative forms, which we shall use in what follows:
\begin{equation}
\label{eq:FS-Lambda}
  \frac{\lambda^x ({-\log\lambda})^{1-s}\Phi(\lambda,s,x)}{\Gamma(1-s)}
  = \sum_{n\in\Z}(1-\Lambda n)^{s-1} e^{2n\pi i x},
  \quad \Lambda=\frac{2\pi i}{\log\lambda},
\end{equation}
and the more compact
\begin{equation}
\label{eq:FS-poles}
  \frac{1}{\Gamma(1-s)}\Phi(\lambda,s,x) = \sum_{a \in S} a^{s-1}e^{a x},
\end{equation}
where $S = \{ a_{n} = 2 \pi i n - \log \lambda : n \in \Z \}$ (this also the set of poles of the generating function of the Apostol-Bernoulli polynomials). It is worth noting that the set $S$ does not depend on which branch of the logarithm is chosen and, therefore, the Lindel\"of-Wirtinger expansion is also independent of this choice, provided that the same branch is used for the term $\lambda^{-x}$ and the $\log \lambda$ in the sum. Indeed, varying the branch is equivalent to shifting the index $n$ in both the pole set $S$ and the expansion.
\end{remark}

In order to obtain asymptotic estimates, we need to order the poles $S$ by their modulus. This ordering varies depending on~$\lambda$. The details are given in the following lemma, which is proved in~\cite{NRV-MC}.


\begin{lemma}
\label{lem:pole-ordering}
Let $a_{n} = 2 \pi i n - \log \lambda$ with $n \in \Z$,
$\lambda \in \C$, $\lambda \ne 0$.
\begin{enumerate}
%
  \item[(a)] If $\Im\lambda > 0$, then for $n \ge 1$, we have
  \[
    0 < |a_{0}| < |a_{1}| < |a_{-1}| < \cdots < |a_{n}| < |a_{-n}| < \cdots
  \]
  \item[(b)] If $\Im\lambda < 0$, then for $n \ge 1$, we have
  \[
    0 < |a_{0}| < |a_{-1}| < |a_{1}| < \cdots < |a_{-n}| < |a_{n}| < \cdots
  \]
  \item[(c)] If $\lambda > 0$, then for $n \ge 1$, we have
  \[
    |a_{0}| < |a_{1}| = |a_{-1}| < \cdots < |a_{n}| = |a_{-n}| < \cdots
  \]
% (note that $a_0 =0$ if and only if $\lambda=1$, in which case it is not a pole,
% hence is excluded in this chain).
  \item[(d)] If $\lambda < 0$, then for $n \ge 1$, we have
  \[
    0 < |a_{0}| = |a_{1}| < |a_{-1}| = |a_{2}| < \cdots < |a_{-n}| = |a_{n+1}| < \cdots
  \]
%
\end{enumerate}
In addition, $|a_{n}| \geq 2\pi (|n| - \frac{1}{2}) \geq \pi$ if $|n| \geq 1$.
\end{lemma}

From here on, we shall only consider partial sums of the series~\eqref{eq:FS-poles},
$\sum_{a\in F}$, over subsets $\emptyset \ne F\subset S$ such that
\begin{equation}
\label{eq:admissible}
  \max\{ |a| : a \in F \} < \min\{ |a| : a \in S \setminus F \}
  \stackrel{\text{\tiny def}}{=} \mu_F.
\end{equation}
This is necessary in order to have the partial sum be of greater order than the tail.
Let us call such subsets \emph{admissible}. The admissible sets vary depending on the cases established in the lemma. Note also that $\mu_{F} > 1$ as mentioned in the last line of the lemma.

The Fourier series of $\lambda^{x} \Phi(\lambda,s,x)$ can be expressed in these terms as follows.

\begin{corollary}
\label{cor:FS-admissible}
Let $\lambda \in \C$ with $0<|\lambda|<1$, $F$ an admissible set in the sense of~\eqref{eq:admissible},
and $s \in \C$ with $\Re s < \alpha < 0$ and $|\Im s| \le \beta$. Then,
\begin{equation*}
  \frac{1}{\Gamma(1-s)}\Phi(\lambda,s,x) = \sum_{a \in F} a^{s-1}e^{a x} + O(\mu_F^{\Re s-1})
\end{equation*}
holds uniformly in $x\in[0,1]$, where the constant implicit in the order term depends only on
$F$, $\alpha$ and $\beta$.
\end{corollary}

\begin{proof}
For simplicity, let us assume that $s \in \R$, with $s<\alpha<0$ (and consequently
$\beta=0$). Taking into account~\eqref{eq:sreal}, the proof of the general case only requires small modifications.

We need to bound the remainder in~\eqref{eq:FS-poles}. Now, if $a\in S\setminus F$,
then $|e^{a x}| = |\lambda^x e^{2\pi i n x}| \le 1$ for $0<|\lambda|<1$, $x \ge 0$, $n\in\Z$,
hence
\[
  \left|\sum_{a\in S\setminus F}a^{s-1}e^{a x}\right| \leq C \mu_F^{s-1}
  \sum_{a\in S\setminus F} \left(\frac{|a|}{\mu_F}\right)^{s-1},
\]
and, since $|a_{n}| \geq 2\pi (|n| - \frac{1}{2}) \geq |n|$, the latter series is bounded by the constant
\begin{equation*}
  C_{\alpha,F} = \sum_{|n|\geq 1} \left(\frac{|n|}{\mu_F}\right)^{\alpha-1} < +\infty.
\qedhere
\end{equation*}
\end{proof}

Using the language of special functions, we have proved that the Fourier series is
a uniformly asymptotic series for $\Phi$ over $x \in [0,1]$, under the stated conditions of $\Re s \to -\infty$ with $\Im s$ bounded.

A priori, the Fourier series need not converge or even be an asymptotic series outside of $[0,1]$, and even if it is, it need not converge or be asymptotic to $\Phi$. We shall prove that in fact the result of Corollary~\ref{cor:FS-admissible} is still valid, uniformly over compact subsets of~$\dom$.


\section{Extension to compact subsets of~$\C$}
\label{sec:extension}

In the case of the Apostol-Bernoulli polynomials, the fact that, in the language of the ``umbral calculus,'' they are a polynomial sequence of binomial type, can be used to prove that the Fourier series is an asymptotic series on arbitrary compact subsets of~$\dom$ (see~\cite{NRV-MC}). This point of view exploits connections between the algebraic and combinatorial properties of polynomial families and their analytic expansions, such as asymptotic series.

Here we adapt the method used in~\cite{NRV-MC} to express the values of
$\Phi(\lambda,s,z)$ in terms of the values at $z=0$, at the cost of having to vary~$s$. Since we are no longer dealing with polynomials, the analytic estimates are more complicated.

\begin{lemma}
\label{lemma:sep}
Let $N\in\N$, $s \in \C$ with $\Re s < 0$ and $\lambda \in \C$ with $0<|\lambda|<1$.
Then, for $z\in\dom$ and $|z|<N$,
\begin{multline*}
  \quad
  \Phi(\lambda,s,z)
  = \sum_{k=0}^{N-1}\frac{\lambda^k}{(k+z)^s}
  + \sum_{n \geq -\Re s-1/2}\binom{-s}{n}z^n\sum_{k=N}^\infty \frac{\lambda^k}{k^{n+s}}
  \\
  + \sum_{n < -\Re s-1/2} \binom{-s}{n} z^n \Biggl(\Phi(\lambda,n+s,0)
  - \sum_{j=1}^{N-1} \frac{\lambda^j}{j^{n+s}}\Biggr).
  \quad
\end{multline*}
\end{lemma}

\begin{proof}
For simplicity, let us assume again that $s \in \R$ and $s < 0$. The general case is taken care of similarly by~\eqref{eq:sreal} (this affects only the justification for changing the order of summation in the infinite series).

We begin by separating the series defining $\Phi$ into a partial sum and a tail:
\[
  \Phi(\lambda,s,z) = \sum_{k=0}^{N-1}\frac{\lambda^k}{(k+z)^s}
  + \sum_{k=N}^\infty\frac{\lambda^k}{(k+z)^s}.
\]
Observe that if $|z|<N$ and $k\geq N$, we may expand the denominators as
\[
  (z+k)^{-s} = k^{-s} \sum_{n=0}^\infty \binom{-s}{n} \frac{z^n}{k^n}
  = \sum_{n<-s-1/2} \binom{-s}{n} \frac{z^n}{k^{n+s}}
  + \sum_{n\geq -s-1/2} \binom{-s}{n} \frac{z^n}{k^{n+s}}.
\]
Then
\begin{align*}
  \sum_{k=N}^\infty\frac{\lambda^k}{(k+z)^s}
  &= \sum_{k=N}^\infty \lambda^k \sum_{n<-s-1/2} \binom{-s}{n} \frac{z^n}{k^{n+s}}
  + \sum_{k=N}^\infty \lambda^k \sum_{n\geq -s-1/2} \binom{-s}{n} \frac{z^n}{k^{n+s}}
  \\
  &= \sum_{n<-s-1/2} \binom{-s}{n}z^n \sum_{k=N}^\infty \frac{\lambda^k}{k^{n+s}}
  + \sum_{n\geq -s-1/2} \binom{-s}{n}z^n \sum_{k=N}^\infty \frac{\lambda^k}{k^{n+s}}.
\end{align*}
The first change in the order of summation is trivial, since one of the series is finite. For the second, we apply Fubini's theorem, noting that if $n\geq -s-1/2$ (which implies $-s\leq n+1$), then $|\binom{-s}{n}|\leq n+1$ and
\begin{multline*}
  \quad
  \sum_{k=N}^\infty |\lambda|^k\,k^{-s}
  \sum_{n\geq -s-1/2} \left|\binom{-s}{n}\right| \frac{|z|^n}{|k|^n}
  \leq \sum_{k=N}^\infty |\lambda|^k\,k^{-s} \sum_{n=0}^\infty (n+1)\frac{|z|^n}{|k|^n}
    \\
  = \sum_{k=N}^\infty |\lambda|^k\,k^{-s} \frac{k^2}{(k-|z|)^2}
  < +\infty,
  \quad
\end{multline*}
since $|\lambda|<1$ (if $s$ is complex and $|\Im s|\leq \beta \in \N$, then an easy bound for $|\binom{-s}{n}|$ is $(n+1)(n+2)\cdots (n+\beta+1)$, and the justification is similar).

Finally, if $s+n<-1/2\ $ ($<0$), then
\[
  \sum_{k=N}^\infty\frac{\lambda^k}{k^{n+s}}
  = \Phi(\lambda,n+s,0) - \sum_{j=1}^{N-1} \frac{\lambda^j}{j^{n+s}}
\]
(since the sum corresponding to $j=0$ is null).
Note that this does not work if $s+n> 0$ since in that case $\Phi(\lambda,n+s,z)$ does not extend continuously to $z=0$.
\end{proof}

We can now state the main theorem.

\begin{theorem}
\label{th:FS-admissible}
Let $\lambda$ and $s$ be complex parameters with $0<|\lambda|<1$ and $\Re s<0$,
$F$~an admissible set in the sense of~\eqref{eq:admissible}, and $K$ a compact set in~$\dom$.
Then, for $z \in K$ and $\Re s\to-\infty$ with $|\Im s|$ bounded,
\begin{equation*}
  \frac{1}{\Gamma(1-s)}\Phi(\lambda,s,z) = \sum_{a \in F} a^{s-1}e^{a z} + O(\mu_F^{\Re s-1}),
\end{equation*}
where the constant implicit in the order term depends only on $F$, $K$, the bound for $|\Im s|$ and~$\lambda$.
\end{theorem}

\begin{proof}
As before, let us assume that $s \in \R$ and $s\to-\infty$. The general case can be proved in a similar way by applying~\eqref{eq:sreal}, and using the complex version of Stirling's Formula, namely
\[
  \Gamma(w) = e^{-w} e^{(w-1/2)\log w} (2\pi)^{1/2} \left(1 + O(w^{-1})\right),
  \ |\arg w| \le \pi-\varepsilon,\ \varepsilon>0,\ |w| \to \infty
\]
to bound the $\Gamma(1-s)$ factors which appear (see, for instance, \cite[\S1.18, p.~47]{Bate}).


Given a compact subset $K$ of $\dom$, fix $N$ such that $|z|< N$ for $z\in K$. Define
\[
  R_F(\lambda,s,z) = \frac{\Phi(\lambda,s,z) }{\Gamma(1-s)}
  - \sum_{a \in F} a^{s-1}e^{a z}.
\]
We shall show that, given a large negative number $s_0$, there exists a constant $C$ depending only on $F$, $K$ and $s_0$, such that
\[
  |R_F(\lambda,s,z)| \leq C \mu_F^{s-1},
  \quad s\leq s_0.
\]
By Lemma~\ref{lemma:sep}, we may separate $R_F(\lambda,s,z)$ into the following four terms:
\begin{align*}
  R_F(\lambda,s,z)
  &= \left(\frac{1}{\Gamma(1-s)} \sum_{k=0}^{N-1}\frac{\lambda^k}{(k+z)^s}\right)
  \\
  &\qquad + \left(\frac{1}{\Gamma(1-s)} \sum_{n\geq -s-1/2}\binom{-s}{n}z^n
  \sum_{k=N}^\infty \frac{\lambda^k}{k^{n+s}}\right)
  \\
  &\qquad - \left(\frac{1}{\Gamma(1-s)} \sum_{n<-s-1/2}\binom{-s}{n}z^n
  \sum_{j=1}^{N-1}\frac{\lambda^j}{j^{n+s}}\right)
  \\
  &\qquad + \left(\frac{1}{\Gamma(1-s)} \sum_{n<-s-1/2} \binom{-s}{n}z^n\Phi(\lambda,n+s,0)
  - \sum_{a \in F} a^{s-1}e^{a z}\right)
  \\
  &= \mathit{(I)} + \mathit{(II)} + \mathit{(III)} + \mathit{(IV)}.
\end{align*}
The finite sum in $\mathit{(I)}$ is easily seen to be uniformly bounded for $z\in K$ by a factor of the form $\gamma^{-s}$ with $\gamma>0$. The presence of $\Gamma(1-s)$ in the denominator implies that $\mathit{(I)}$ is actually $O(\delta^{s-1})$ for \emph{any} $\delta>1$, not just for $\delta=\mu_F$.


For the sum that appears in $\mathit{(II)}$, since $|\binom{-s}{n}| \leq n+1$ and $|z| \leq r < N$ for $z\in K$, with $r$ depending only on the compact set $K$, we get the bound
\[
  \sum_{k=N}^\infty \frac{|\lambda|^k}{k^s} \sum_{n\geq -s-1/2}   (n+1)\frac{r^n}{k^n}.
\]
Noting that, for $b<1$ and $m$ an integer,
\[
  \sum_{n=m}^\infty (n+1)b^n = \frac{b^{m}}{(1-b)^2}(1+m(1-b)),
\]
and since $\sum_{n\geq -s-1/2} = \sum_{n=m}^\infty$ with $-s-2\leq m\leq -s+1$, we have, for any $k\geq N$, the bound
\[
  \sum_{n\geq -s-1/2}^{} (n+1)\frac{r^n}{k^n}
  \leq \left(\frac{r}{k}\right)^{-s-2}\frac{1+(-s+1)(1-r/k)}{(1-r/k)^2}
  \leq \frac{k^4N^{-s-2}(2-s)}{k^{-s}(N-r)}.
\]
Thus the sum in $\mathit{(II)}$ is bounded for
\[
  \frac{N^{-s-2}(2-s)}{N-r} \sum_{k=N}^\infty k^4|\lambda|^k = C(2-s)N^{-s-1}
\]
where the constant $C$ depends only on $N$, $N-r$ and $\lambda$, and hence only on the compact set $K$ and~$\lambda$.
Again, the $\Gamma(1-s)$ term in the denominator shows that $\mathit{(II)}$ has order
$O(\delta^{s-1})$ for any $\delta>1$.

In $\mathit{(III)}$, since the sum in $j$ is finite, we consider each such summand separately and
study the behavior as $s \to -\infty$ of the expressions
\[
  \frac{1}{\Gamma(1-s)} \sum_{n<-s-1/2} \binom{-s}{n} \frac{|z|^n}{a^{n+s}},
\]
where $a=1,2,\dots, N-1$ and $|z|<N$.
Note that $a^n\geq 1$. Dividing and multiplying by $N^{-s}$ and noting that
$0 < \binom{-s}{n} < 2\cdot 2^{-s}$, we get the following bound for the modulus of~$\mathit{(III)}$:
\begin{align*}
  \frac{N^{-s}}{a^{s}\Gamma(1-s)}
  & \sum_{n<-s-1/2}\binom{-s}{n}\frac{|z|^n}{N^{-s}}
  \leq \frac{N^{-s}}{a^{s}\Gamma(1-s)} \sum_{n<-s-1/2}\binom{-s}{n} \frac{|z|^n}{N^{n+1/2}}
  \\
  &\leq 2 \, \frac{N^{-s-1/2}2^{-s}}{a^{s}\Gamma(1-s)} \left(1-\frac{|z|}{N}\right)^{-1}
  \leq  C_K \, \frac{N^{-s-1/2}2^{-s}}{a^{s}\Gamma(1-s)}.
\end{align*}
Once again, the $\Gamma(1-s)$ term in the denominator makes $\mathit{(III)}$ of the order $O(\delta^{s-1})$ for any $\delta>1$.

The only subtlety in obtaining the bound stated above lies with~$\mathit{(IV)}$.
Observe that $n+s<-1/2<0$. Then, by Corollary~\ref{cor:FS-admissible} with $\alpha=-1/2$ and $x=0$, there is a constant $C$ depending only on $F$ such that
\begin{equation}
\label{eq:asintoticaen0}
  |R_F(\lambda,n+s,0)| = \left|\frac{\Phi(\lambda,n+s,0)}{\Gamma(1-n-s)}
  - \sum_{a \in F} a^{s+n-1}\right| \leq C \mu_F^{{n+}s-1}.
\end{equation}
Substituting
\[
  \Phi(\lambda,n+s,0) = \Gamma(1-n-s) \biggl(R_F(\lambda,n+s,0)+\sum_{a \in F} a^{s+n-1}\biggr)
\]
into $\mathit{(IV)}$ and noting that
\[
  \binom{-s}{n} \frac{\Gamma(1-s-n)}{\Gamma(1-s)}=\frac{1}{n!},
\]
we see that $\mathit{(IV)}$ may be rewritten as
\begin{gather} \nonumber
  \sum_{n<-s-1/2} \frac{z^n}{n!} \biggl( R_F(\lambda,n+s,0)
  + \sum_{a \in F} a^{s+n-1} \biggr) - \sum_{a \in F} a^{s-1}e^{az}
  \\
  \label{eq:IVrewritten}
  = \sum_{n<-s-1/2} \frac{z^n}{n!} R_F(\lambda,n+s,0) +
  \sum_{a \in F}a^{s-1} \left( \sum_{n<-s-1/2}\frac{(az)^n}{n!}-e^{az} \right).
\end{gather}
By~\eqref{eq:asintoticaen0}, the first of these summands in~\eqref{eq:IVrewritten} is bounded by
\[
  C\sum_{n<-s-1/2} \frac{|z|^n}{n!} \mu_F^{n+s-1} \leq C\mu_F^{s-1}e^{|z|\mu_F}
  \leq C_{K,F}\mu_F^{s-1}.
\]

For the second summand in~\eqref{eq:IVrewritten}, we note that, for any $z\in K$ and any $a\in F$, there is a constant $C=C_{K,F}$ such that
\[
  \left| \sum_{n<-s-1/2}\frac{(az)^n}{n!}-e^{az} \right|
  \leq \sum_{n\geq -s-1/2}\frac{|az|^n}{n!}\leq \sum_{n\geq -s-1/2}\frac{C^n}{n!}.
\]
To bound this tail of the exponential series, recall that $\sum_{n\geq -s-1/2} = \sum_{n=m}^\infty$ with $-s-2\leq m\leq -s+1$. For the sum over $n \geq�m$, the Lagrange form of the remainder for Taylor series gives a bound of the form $e^{C} {C^{m}}/{m!}$, from which the bound $C^{-s+1}/\Gamma(-s-2)$, for a suitable new constant, now follows.
This is once again of order $O(\delta^{s-1})$ for any $\delta>1$; thus, in the entire process of bounding the different terms in $\mathit{(I)} + \mathit{(II)} + \mathit{(III)} + \mathit{(IV)}$, the limitation given by $\mu_F^{s-1}$ is due exclusively to the first term in~\eqref{eq:IVrewritten}.
\end{proof}

\section{Asymptotic behavior of $\Phi(\lambda,s,z)$ on $\C$ when $s\to-\infty$}
\label{sec:asymptotic}

We note here two simple consequences of our main theorem.

\begin{corollary}
\label{cor:lambdanonneg}
Let $0<|\lambda|<1$ and $\lambda \notin ({-1},0)$. Then,
\begin{equation}
\label{eq:lambdanonneg}
  \lim_{\R\ni s\to-\infty} \frac{({-\log\lambda})^{1-s}\Phi(\lambda,s,z)}{\Gamma(1-s)} = \lambda^{-z}
\end{equation}
uniformly for $z$ on compact subsets of~$\dom$.
\end{corollary}

\begin{proof}
The restriction on $\lambda$ implies that $F = \{a_0=-\log \lambda\}$ is an admissible set.
Then, by Theorem~\ref{th:FS-admissible},
\begin{equation*}
  \frac{\Phi(\lambda,s,z)}{\Gamma(1-s)} = ({-\log\lambda})^{s-1}\lambda^{-z}+ O(\mu_F^{\Re s-1}).
\end{equation*}
Dividing both sides of this equation by $({-\log\lambda})^{s-1}$ and noting that $\mu_F/|{-\log\lambda}|>1$, we obtain the result.
\end{proof}

\begin{corollary}
\label{cor:lambdaneg}
Let $\lambda \in ({-1},0)$. Then,
\begin{equation}
\label{eq:lambdaneg}
  \lim_{\R\ni s\to-\infty} \left(\frac{(-\log|\lambda|-i\pi)^{1-s}
  \Phi(\lambda,s,z)}{\Gamma(1-s)}\lambda^{z}
  - \frac{(-\log|\lambda|-i\pi)^{1-s}}{(-\log|\lambda|+i\pi)^{1-s}}e^{2\pi iz} \right) = 1
\end{equation}
uniformly for $z$ on compact subsets of~$\dom$.
\end{corollary}

\begin{proof}
In this case, $F = \{a_0=-\log\lambda,\,a_1=2\pi i-\log\lambda\}$ is an admissible set
(note that $|{-\log\lambda}| = |2\pi i-\log\lambda|$ and therefore there is no admissible set with a single point). By Theorem~\ref{th:FS-admissible}, we have
\begin{equation*}
  \frac{\Phi(\lambda,s,z)}{\Gamma(1-s)}
  = ({-\log \lambda})^{s-1}\lambda^{-z} + (2\pi i-\log\lambda)^{s-1}\lambda^{-z}e^{2\pi iz}
  + O(\mu_F^{\Re s-1}),
\end{equation*}
and some simple manipulations lead to the result.
\end{proof}


Observe that the term $\frac{\log|\lambda|+i\pi}{\log|\lambda|-i\pi}$ in the last corollary has modulus one and therefore the Lerch function exhibits oscillatory behavior. Let us examine a simple case.
With $\lambda=-e^{-\pi}$, we have
\begin{equation*}
  \frac{\log|\lambda|+i\pi}{\log|\lambda|-i\pi}=-i.
\end{equation*}
Fixing $\alpha\in[0,1)$ and setting $s=1-n-\alpha$, $n\in\N$,
for $n$ a multiple of~$4$, Corollary~\ref{cor:lambdaneg} says that
\begin{equation*}
  \frac{({-1})^\alpha e^{\alpha i\pi/4}(-\pi+i\pi)^{n+\alpha}\Phi(\lambda,n+\alpha,z)} {\Gamma(n+\alpha)}
  \longrightarrow 2e^{\pi z} \cos\left(\pi z-\frac{\alpha\pi}{4}\right),
  \quad n\to \infty,
\end{equation*}
while if $n$ is congruent to $2$ modulo $4$, we have
\begin{equation*}
  \frac{({-1})^\alpha e^{\alpha i\pi/4}({-\pi}+i\pi)^{n+\alpha}\Phi(\lambda,n+\alpha,z)} {\Gamma(n+\alpha)}
  \longrightarrow 2ie^{\pi z} \sin\left(\pi z-\frac{\alpha\pi}{4}\right),
\quad n\to \infty.
\end{equation*}
For $\alpha=0$, we obtain the Apostol-Bernoulli polynomials, whose oscillatory behavior has been studied in~\cite{NRV-MC}.

\begin{remark}
For negative $\lambda$, we have the straightforward relation
\begin{equation*}
  \Phi(\lambda,s,z) = 2^{-s} \left(\Phi\Bigl(\lambda^2,s,\frac{z}{2}\Bigr)
  + \lambda \Phi\Bigl(\lambda^2,s,\frac{z+1}{2}\Bigr)\right),
\end{equation*}
which implies that
\begin{multline*}
  \frac{({-\log\lambda})^{1-s}\Phi(\lambda,s,z)}{\Gamma(1-s)}
  \\
  = \frac12 \frac{({-\log\lambda^2})^{1-s}\Phi(\lambda^2,s,\frac{z}{2})}{\Gamma(1-s)}
  + \frac12 \lambda \frac{({-\log\lambda^2})^{1-s} \Phi(\lambda^2,s,\frac{z+1}{2})}{\Gamma(1-s)}.
\end{multline*}
However, if we apply Corollary~\ref{cor:lambdanonneg}, then we only obtain
\begin{equation*}
  \lim_{\R\ni s\to-\infty} \frac{({-\log\lambda})^{1-s}\Phi(\lambda,s,z)}{\Gamma(1-s)}
  = \frac12\left(|\lambda|^{-z}+ \lambda |\lambda|^{-z-1}\right) = 0
\end{equation*}
which does not yield the asymptotic behavior of the Lerch function.
\end{remark}


\section{A nice application: polylogarithms}
\label{sec:polylog}

Recall that the polylogarithm functions are defined by the series
\[
  \Li_{s}(\lambda) = \sum_{k=1}^{\infty} \frac{\lambda^{k}}{k^{s}},
\]
and thus $\Li_{s}(\lambda) = \lambda \Phi(\lambda,s,1)$. For $\Re s < 0$, we also have $\Li_{s}(\lambda) = \Phi(\lambda,s,0)$.

Setting $z = 0$ in Corollaries~\ref{cor:lambdanonneg} and~\ref{cor:lambdaneg},
and letting $\lambda = e^{\mu}$ with $\Re \mu < 0$ and $\Im \mu \in (-\pi,\pi]$
(the boundary case $\Im \mu = \pi$ now corresponds to Corollary~\ref{cor:lambdaneg}),
we have $\log \lambda = \mu$ and hence
\begin{equation*}
%\label{eq:polylog1}
  \lim_{\R \ni s \to -\infty}
  \frac{(-\mu)^{1-s} \Li_{s}(e^{\mu})}{\Gamma(1-s)} = 1,
  \qquad
  \Im \mu \in (-\pi,\pi),
\end{equation*}
and
\begin{equation*}
%\label{eq:polylog2}
  \lim_{\R \ni s \to -\infty}
  \left(
      \frac{(-\mu)^{1-s} \Li_{s}(e^{\mu})}{\Gamma(1-s)}
    - \frac{(-\mu)^{1-s}}{(-\mu + 2 \pi i)^{1-s}}
  \right)
    = 1,
  \qquad
  \Im \mu = \pi.
\end{equation*}
Thus one also obtains the asymptotic behavior of the polylogarithm functions in the stated cases.

\subsection*{Acknowledgement.}
We wish to thank the referee for the careful reading of the paper and for suggesting various corrections and improvements.


%--------------------------------------
\begin{thebibliography}{00}
%--------------------------------------

\bibitem{Ap}
T. M. Apostol,
On the Lerch zeta function,
\textit{Pacific J. Math.}
\textbf{1} (1951), 161--167;
addendum, \textit{Pacific J. Math.} \textbf{2} (1952), 10.

\bibitem{Ap-NIST}
T. M. Apostol,
Zeta and related functions,
\textit{NIST handbook of mathematical functions}
(edited by F.~W. F.~Olver, D.~W. Lozier, R.~F. Boisvert and C.~W. Clark),
601--616,
National Institute of Standards and Technology, Washington, DC, and
Cambridge University Press, Cambridge, 2010.
Available online in \url{http://dlmf.nist.gov/25}

\bibitem{BaSa}
E. P. Balanzario and J. S\'{a}nchez-Ortiz,
Riemann-Siegel integral formula for the Lerch zeta function,
\textit{Math. Comp.},
to appear.
DOI: \url{http://dx.doi.org/10.1090/S0025-5718-2011-02566-4}

\bibitem{Boy}
K. N. Boyadzhiev,
Apostol-Bernoulli functions, derivative polynomials and Eulerian polynomials,
\textit{Adv. Appl. Discrete Math.}
\textbf{1} (2008), 109--122.

% \cite[\S1.11, p.~27]{Bate}
\bibitem{Bate} % Bateman project
A. Erd\'{e}lyi, W. Magnus, F. Oberhettinger and F. Tricomi,
\textit{Higher Transcendental Functions}, Volume I,
McGraw-Hill, New York, 1953.

%% http://dlmf.nist.gov/25.14
%% (el cap\'itulo 25, ``Zeta and related functions'', es de Tom M. Apostol)
%% \cite[\S25.14]{DLMF-online}
%\bibitem{DLMF-online}
%NIST,
%\textit{Digital Library of Mathematical Functions},
%National Institute of Standards and Technology, 2010, % Release date 2010-05-07
%\url{http://dlmf.nist.gov/}

%%% El anterior en papel:
%\bibitem{DLMF-printed}
%\textit{NIST Handbook of Mathematical Functions}
%(edited by F. W. F. Olver, D. W. Lozier, R. F. Boisvert and C. W. Clark),
%National Institute of Standards and Technology and
%Cambridge University Press, 2010.

\bibitem{FL}
C. Ferreira and J. L. L\'{o}pez,
Asymptotic expansions of the Hurwitz-Lerch zeta function,
\textit{J. Math. Anal. Appl.}
\textbf{298} (2004), no. 1, 210--224.

\bibitem{GuiS}
J. Guillera and J. Sondow,
Double integrals and infinite products for some classical constants via analytic continuations of Lerch's transcendent,
\textit{Ramanujan J.}
\textbf{16} (2008), 247--270.
% Preprint in \url{http://arxiv.org/abs/math.NT/0506319}

\bibitem{Le}
M. Lerch,
Note sur la fonction $\mathfrak{K}(w,x,s) = \sum_{k=0}^{\infty} e^{2\pi ikx}(w+k)^{-s}$,
\textit{Acta Math.}
\textbf{11} (1887), 19--24.
% \url{http://www.springerlink.com/content/bpn8104100w31667/}
% http://www.springerlink.com/content/0001-5962/11/1-4/

\bibitem{MOS}
W. Magnus, F. Oberhettinger and R. P. Soni,
\textit{Formulas and Theorems for the Special Functions of Mathematical Physics},
3rd ed., Springer, New York, 1966.

\bibitem{NRV-JAT}
L. M. Navas, F. J. Ruiz and J. L. Varona,
The M\"obius inversion formula for Fourier series
applied to Bernoulli and Euler polynomials,
\textit{J. Approx. Theory}
\textbf{163} (2011), 22--40.

\bibitem{NRV-MC}
L. M. Navas, F. J. Ruiz and J. L. Varona,
Asymptotic estimates for Apostol-Bernoulli and Apostol-Euler polynomials,
\textit{Math. Comp.},
to appear.
DOI: \url{http://dx.doi.org/10.1090/S0025-5718-2012-02568-3}
% Preprint in \url{http://arxiv.org/abs/1102.1493}

\bibitem{SC-book}
H. M. Srivastava and J. Choi,
\textit{Series associated with the zeta and related functions},
Kluwer Academic, Dordrecht, 2001.

\bibitem{Wir}
W. Wirtinger,
\"Uber eine besondere Dirichletsche Reihe,
\textit{J. Reine Angew. Math.} \textbf{129} (1905), 214--219.
% % Journal f\"ur die Reine und Angewandte Mathematik

%--------------------------------------
\end{thebibliography}
%--------------------------------------

%--------------------------------------
\end{document}
%--------------------------------------
